\section{Experimental Evaluation Details}\label{sec:app_experimental_evaluation}


\subsection{Frozen Subsets}\label{sec:app_bench_lists}

The frozen subsets \texttt{FC100SolvedSet1} and \texttt{FC100OpenSet1} consist of 100 problems each, sampled uniformly at random from statements in the \texttt{research solved} and \texttt{research open} categories, respectively.
They are defined by the files \path{FormalConjectures/Subsets/FC100OpenSet1.lean} and \path{FormalConjectures/Subsets/FC100SolvedSet1.lean} in the repository, which import exactly the corresponding theorem statements.
These files are compiled as part of every tagged release, ensuring that the problem sets remain well-defined and their statements compile correctly across all supported Lean versions.

The subset names (e.g., \texttt{FC100OpenSet1}) are \emph{orthogonal} to the repository's release tags (\texttt{bench-v$N$-lean4.$X$.$Y$}, see Section~\ref{sec:versioning}).
When the bench version $N$ stays the same and only the Lean version changes (e.g., from \texttt{bench-v1-lean4.27.0} to \texttt{bench-v1-lean4.29.0}), the problem set is guaranteed to contain the same statements; the new tag only ensures compilation against the updated Lean toolchain and Mathlib.
When the bench version increments (e.g., from \texttt{bench-v1} to \texttt{bench-v2}), statements in an existing set may receive fixes (e.g., corrected misformalizations) or category updates (e.g., an \texttt{open} problem reclassified as \texttt{solved}), but the \emph{membership} of the set remains stable: no problems are added or removed.
For full reproducibility, results should specify both the set name and the release tag, e.g., \texttt{FC100SolvedSet1@bench-v1-lean4.27.0}.

In the future, we plan to release additional sets (e.g., \texttt{FC100SolvedSet2}, \texttt{FC500OpenSet1}) that may contain entirely different problems, enabling targeted evaluations at different scales or difficulty levels while preserving earlier sets.


\subsection{Illustrative Evaluation Details}\label{sec:illustrative_evaluation}


As noted in Section \ref{sec:experimental_results}, our evaluations focus on demonstrating the benchmark’s utility and its sensitivity to scaling.
For all systems, a solution is correct if and only if the Lean~4 kernel accepts the proof term without relying on forbidden axioms.


\paragraph{AlphaProof Setup.}
We evaluated a slightly updated version of AlphaProof \citep{AlphaProof25}, with tree-search inference, without the test-time reinforcement learning loop.
As reported in Table \ref{tab:results_bench_table}, we distinguish between two compute configurations for AlphaProof to demonstrate the benchmark's sensitivity to search budget: 1k sims utilizes 1,000 simulations and 2 attempts per problem (approximately 0.1 TPUh on v6e TPUs); and 16k sims and 2 attempts (approximately 1.6 TPUh on v6e TPUs).


\paragraph{DeepMind Prover Agent Setup.}
The DeepMind prover agent \citep{DMProverAgent26} results were obtained using a development version of the method under active experimentation.
Based on current API pricing, this evaluation costs up to \$100 per problem.

